% !TeX root = ../../Book4.tex
% 纲要标题：### 9.5 维数平移与合冲模初步
% 纲要标签：b4:sec:0805
% 纲要位置：计划纲要.md，第 3426 行。

\section{维数平移与合冲模初步}
\label{b4:sec:0805}

短正合列中的投射或内射项消去高次导出群，使相邻次数之间出现自然同构。逐次取合冲模或余合冲模，便把高次计算移到更低次数。


% 纲要内容（仅作写作计划，尚非教材正文）：
%
% * 对短正合列 \(0\to\Omega M\to P\to M\to0\) 及投射模 \(P\)，由长正合列推出正次数 Ext 和 Tor 的维数平移，并注明变量侧别
% * 对内射嵌入 \(0\to N\to I\to C\to0\) 给出对偶的 Ext 平移公式；区分合冲模选择与导出不变量的自然同构
% * 迭代平移与截短消解，说明在有限步后出现投射合冲模的含义；不在此预设一般同调维数理论
% * 本节是第三卷第15–21章所用同调工具包的一部分；第11.1–11.2节在其基础上统一定义和研究各类同调维数
%
% ---
%

\subsection{投射消解中的 Ext 与 Tor 维数平移}
\label{b4:subsec:0805:01}

\begin{definition}\label{b4:def:syzygy-basic}
设 \(R\) 是含幺环，\(M\) 是幺性左 \(R\) 模。给定投射满射 \(P\to M\)，其核记为 \(\Omega M\)，称为与这一选择对应的第一\term[hechongmo]{合冲模}[syzygy module]。在固定投射消解 \(P_\bullet\xrightarrow{\varepsilon}M\) 中，令 \(\Omega^0M=M\)、\(\pi_0=\varepsilon:P_0\to M\)，并递归定义
\[
\Omega^{r+1}M=\ker\pi_r,\qquad \pi_r:P_r\twoheadrightarrow\Omega^rM.
\]
这里 \(r\ge1\) 时，消解的正合性给出 \(\operatorname{im}\mathrm{d}_r=\Omega^rM\subseteq P_{r-1}\)，所以微分 \(\mathrm{d}_r:P_r\to P_{r-1}\) 唯一因子化为满射 \(\pi_r\) 与包含映射的复合。这就确定了递归中使用的映射。右模用右投射消解作同一定义。
\end{definition}
\symidx{\(\Omega^rM\)：相对于选定投射消解的第 \(r\) 次合冲模}

\begin{theorem}[基础维数平移]\label{b4:thm:dimension-shift-basic}
设 \(R\) 是含幺环，\(0\to K\to P\to M\to0\) 是幺性左 \(R\) 模短正合列，且 \(P\) 投射。则对任意幺性左 \(R\) 模 \(N\)、幺性右 \(R\) 模 \(A\) 及整数 \(i\ge1\)，连接态射给出同构
\[
\operatorname{Ext}_R^i(K,N)\cong\operatorname{Ext}_R^{i+1}(M,N),
\qquad
\operatorname{Tor}_{i+1}^R(A,M)\cong\operatorname{Tor}_i^R(A,K).
\]
Ext 同构对 \(N\) 及所给短正合列之间的态射自然；Tor 同构对 \(A\) 及所给短正合列之间的态射自然。低次则为正合列
\[
\begin{aligned}
0&\longrightarrow\operatorname{Hom}_R(M,N)\longrightarrow\operatorname{Hom}_R(P,N)
\longrightarrow\operatorname{Hom}_R(K,N)\\
&\longrightarrow\operatorname{Ext}_R^1(M,N)\longrightarrow0,\\
0&\longrightarrow\operatorname{Tor}_1^R(A,M)\longrightarrow A\otimes_RK
\longrightarrow A\otimes_RP\\
&\longrightarrow A\otimes_RM\longrightarrow0.
\end{aligned}
\]
若原列为幺性右 \(R\) 模短正合列，则 Tor 结论改为 \(\operatorname{Tor}_{i+1}^R(M,N)\cong\operatorname{Tor}_i^R(K,N)\)，其中 \(N\) 为幺性左 \(R\) 模；该同构对 \(N\) 及所给右模短正合列之间的态射自然。
\end{theorem}
\begin{proof}
第一变量 Ext 长正合列在相关位置为
\[
\operatorname{Ext}_R^i(P,N)\longrightarrow\operatorname{Ext}_R^i(K,N)
\longrightarrow\operatorname{Ext}_R^{i+1}(M,N)
\longrightarrow\operatorname{Ext}_R^{i+1}(P,N).
\]
\(i\ge1\) 时两端由投射性消失。Tor 的相应片段是
\[
\operatorname{Tor}_{i+1}^R(A,P)\longrightarrow\operatorname{Tor}_{i+1}^R(A,M)
\longrightarrow\operatorname{Tor}_i^R(A,K)\longrightarrow\operatorname{Tor}_i^R(A,P),
\]
两端由投射蕴含平坦而消失。当 \(i=0\) 时，Ext 片段左端是 \(\operatorname{Hom}_R(P,N)\)，Tor 片段右端是 \(A\otimes_RP\)，它们一般不为零，因此连接态射未必是同构；保留这些零次项才得到所列低次正合列。自然性来自原长正合列对短正合列态射的自然性；右模情形把同一证明用于另一变量。
\end{proof}

这一操作称为\term[weishu-pingyi]{维数平移}[dimension shifting]。一条短正合列只使次数降低一步；逐次取合冲模并重复这一操作，才能把高次问题降到一次。零次仍须使用包含 Hom 或张量项的正合列。

\subsection{内射消解中的 Ext 平移与自然性}
\label{b4:subsec:0805:02}

设 \(R\) 是含幺环，\(N\) 是幺性左 \(R\) 模。给定到内射模的单射 \(j:N\to I\)，其余核 \(C=I/j(N)\) 称为与这一选择对应的第一\term[yuhechongmo]{余合冲模}[cosyzygy module]。在内射消解 \(0\to N\to I^0\xrightarrow{\mathrm{d}^0}I^1\xrightarrow{\mathrm{d}^1}\cdots\) 中，第一次取商给出
\[
C^1=I^0/N\cong\operatorname{im}\mathrm{d}^0=\ker \mathrm{d}^1\subseteq I^1.
\]
因此 \(\mathrm{d}^0\) 在这个商上诱导单射 \(C^1\to I^1\)，再取 \(C^2=I^1/C^1\)，就有短正合列 \(0\to C^1\to I^1\to C^2\to0\)。这里将 \(N\) 与其在 \(I^0\) 中的像识别，并将 \(C^1\) 与上述单射的像识别。不断取这样的余核，与投射消解中不断取核相对偶。

\begin{proposition}\label{b4:prop:injective-dimension-shift}
设 \(R\) 是含幺环，\(0\to N\to I\to C\to0\) 为幺性左 \(R\) 模短正合列，\(I\) 内射。对每个幺性左 \(R\) 模 \(M\) 和整数 \(i\ge1\)，有同构
\[
\operatorname{Ext}_R^i(M,C)\cong\operatorname{Ext}_R^{i+1}(M,N).
\]
该同构对 \(M\) 及所给短正合列之间的态射自然。零次给出
\[
\operatorname{Ext}_R^1(M,N)\cong
\operatorname{coker}\bigl(\operatorname{Hom}_R(M,I)\to\operatorname{Hom}_R(M,C)\bigr).
\]
\end{proposition}
\begin{proof}
第二变量长正合列在正次的片段为
\[
\operatorname{Ext}_R^i(M,I)\longrightarrow\operatorname{Ext}_R^i(M,C)
\longrightarrow\operatorname{Ext}_R^{i+1}(M,N)
\longrightarrow\operatorname{Ext}_R^{i+1}(M,I).
\]
\(i\ge1\) 时两端由内射性消失，故中间连接态射为同构；零次则留下
\[
\begin{aligned}
0&\longrightarrow\operatorname{Hom}_R(M,N)\longrightarrow\operatorname{Hom}_R(M,I)
\longrightarrow\operatorname{Hom}_R(M,C)\\
&\longrightarrow\operatorname{Ext}_R^1(M,N)\longrightarrow0,
\end{aligned}
\]
从而得到所述余核表达。连接态射的自然性来自长正合列的自然性。
\end{proof}

余合冲模 \(C\) 取决于所选内射嵌入，正如 \(\Omega M\) 取决于所选投射满射。这里的自然性所比较的是给定短正合列及其态射；更换消解时，比较定理提供相应的比较映射。

\begin{lemma}[Schanuel 引理]\label{b4:lem:schanuel}
设 \(R\) 是含幺环，\(0\to K\to P\to M\to0\)、\(0\to K'\to P'\to M\to0\) 为两个幺性左 \(R\) 模短正合列。若中间模 \(P,P'\) 都投射，则 \(K\oplus P'\cong K'\oplus P\)。
\end{lemma}
\begin{proof}
记两个满射为 \(\pi:P\to M\)、\(\pi':P'\to M\)，取纤维积
\[
E=P\times_M P'=\{(p,p')\in P\oplus P':\pi(p)=\pi'(p')\}.
\]
任取 \(p'\in P'\)，由 \(\pi\) 满射，可取 \(p\in P\) 使 \(\pi(p)=\pi'(p')\)，于是 \((p,p')\in E\) 映到 \(p'\)。故第二投影 \(E\to P'\) 满射，其核为 \(\{(p,0):\pi(p)=0\}\cong K\)；同理第一投影 \(E\to P\) 满射，其核为 \(\{(0,p'):\pi'(p')=0\}\cong K'\)。这样得到短正合列
\[
0\longrightarrow K\longrightarrow E\longrightarrow P'\longrightarrow0,
\qquad
0\longrightarrow K'\longrightarrow E\longrightarrow P\longrightarrow0.
\]
因右端投射，两列均分裂，故 \(E\cong K\oplus P'\cong K'\oplus P\)。
\end{proof}

这个\term[schanuel-yinli]{Schanuel 引理}[Schanuel's lemma]说明两种合冲模在加上投射直和因子后同构。\footnote{沙努埃尔（Stephen Hoel Schanuel，1933-2014），美国数学家，研究同调代数、范畴论与数论。}这不能推出 \(K\cong K'\)，两者甚至可能不同构。例如在非零环 \(R\) 上，自由模 \(M\) 可以取恒等满射，核为零；也可以取投影 \(M\oplus R\to M\)，核为非零的 \(R\)。此外，两条短正合列的分裂一般不唯一，所得直和同构也依赖于分裂的选择。

\subsection{迭代平移与截短消解}
\label{b4:subsec:0805:03}

\begin{proposition}\label{b4:prop:iterated-dimension-shift}
设 \(R\) 是含幺环，\(M,N\) 是幺性左 \(R\) 模，\(A\) 是幺性右 \(R\) 模。固定 \(M\) 的投射消解 \(P_\bullet\to M\) 及其 \(r\) 次合冲模 \(\Omega^rM\)，其中 \(r\ge0\) 为整数。则对整数 \(i\ge1\)，有
\[
\operatorname{Ext}_R^i(\Omega^rM,N)\cong\operatorname{Ext}_R^{i+r}(M,N),
\qquad
\operatorname{Tor}_i^R(A,\Omega^rM)\cong\operatorname{Tor}_{i+r}^R(A,M).
\]
若固定 \(N\) 的内射消解 \(N\to I^\bullet\)，令 \(C^0=N\)、\(C^{j+1}=\operatorname{coker}(C^j\to I^j)\)（\(j\ge0\)）。这里 \(C^0\to I^0\) 为增广单射；对 \(j\ge1\)，内射消解的正合性使 \(I^{j-1}\to I^j\) 在前一步的商 \(C^j\) 上诱导单射，其像为 \(\ker(I^j\to I^{j+1})\)。则
\[
\operatorname{Ext}_R^i(M,C^r)\cong\operatorname{Ext}_R^{i+r}(M,N).
\]
投射平移的两个同构对另一变量及所给投射消解之间的增广链映射自然；内射平移的同构对 \(M\) 及所给内射消解之间的增广上链映射自然。此外，固定消解中的 \(\Omega^rM\) 投射，当且仅当对所有幺性左 \(R\) 模 \(N\)，\(\operatorname{Ext}_R^{r+1}(M,N)=0\)。
\end{proposition}
\begin{proof}
对各短正合列 \(0\to\Omega^{j+1}M\to P_j\to\Omega^jM\to0\) 依次用\cref{b4:thm:dimension-shift-basic}。例如两步给出
\[
\operatorname{Ext}_R^{i+2}(M,N)\cong\operatorname{Ext}_R^{i+1}(\Omega M,N)
\cong\operatorname{Ext}_R^i(\Omega^2M,N).
\]
每次将次数降低一，同时将第一变量换成下一次合冲模；复合 \(r\) 次连接同构即得第一式。Tor 用相同的短正合列迭代，内射式则对 \(0\to C^j\to I^j\to C^{j+1}\to0\) 反复用\cref{b4:prop:injective-dimension-shift}。增广链映射在各核上诱导映射，增广上链映射在各余核上诱导映射；这些映射组成短正合列之间的态射，故每一步的连接同构及其复合都自然。

最后，对每个幺性左 \(R\) 模 \(N\)，第一式在 \(i=1\) 时成为
\[
\operatorname{Ext}_R^1(\Omega^rM,N)\cong\operatorname{Ext}_R^{r+1}(M,N).
\]
因此右边对所有这样的 \(N\) 消失，当且仅当左边也对所有 \(N\) 消失，而\cref{b4:thm:ext-projective-injective}说明后者当且仅当 \(\Omega^rM\) 投射。
\end{proof}

若 \(r=0\)，这一条件就是 \(M\) 本身投射。若 \(r\ge1\) 且 \(\Omega^rM\) 投射，原消解可截成
\[
0\longrightarrow\Omega^rM\longrightarrow P_{r-1}\longrightarrow\cdots
\longrightarrow P_0\longrightarrow M\longrightarrow0.
\]
这里把 \(\Omega^rM\) 作为第 \(r\) 次投射项，所谓长度不超过 \(r\)，指投射项的最高非零次数不超过 \(r\)，不把增广目标 \(M\) 计作投射项。例如长度零表示 \(M\) 投射；长度不超过一表示存在 \(0\to P_1\to P_0\to M\to0\)，其中 \(P_0,P_1\) 均投射。反过来，若有长度不超过 \(r\) 的投射消解，施加 \(\operatorname{Hom}_R(-,N)\) 后，第 \(r+1\) 次项为零，所以对每个 \(N\) 都有 \(\operatorname{Ext}_R^{r+1}(M,N)=0\)。由上述判据，任意投射消解的第 \(r\) 次合冲模都投射。高次 Ext 的全称消失与投射消解能在第 \(r\) 次截短因此等价；\chapref{b4:ch:10}再用它定义并比较各种同调维数。

\subsection{维数平移在局部代数中的应用}
\label{b4:subsec:0805:04}

第三卷的局部代数以各项有限秩的自由消解记录生成元和关系。这里已经得到其中所需的基本同调推理：把 \(0\to\Omega M\to P\to M\to0\) 送入 Ext 或 Tor 的长正合列，利用 \(P\) 的正次导出函子消失，把次数降低一步。有限性和局部性在选择极小自由消解、利用 Nakayama 引理\footnote{中山正（假名：なかやま ただし，罗马音：Tadashi Nakayama，1912-1964），日本数学家，研究环论与表示论。}时另有作用；上述维数平移本身对任意含幺环成立。

\begin{example}\label{b4:exm:local-shift-koszul}
设 \(k\) 是域，\(R=k[x,y]_{(x,y)}\)，\(\mathfrak m=(x,y)R\)，剩余域记为 \(k\)。由于 \(R\) 交换，以下各左模也通过 \(m\cdot r=r\cdot m\) 视为右 \(R\) 模。\(R\) 是整环，\(R/(x)\cong k[y]_{(y)}\) 也是整环，所以 \(x,y\) 满足\cref{b4:prop:koszul-resolution}的假设，给出自由消解
\[
0\longrightarrow R\xrightarrow{\binom{-y}{x}}R^2
\xrightarrow{(x\;y)}R\longrightarrow k\longrightarrow0.
\]
与 \(k\) 张量后，两个微分映射均为零，故 \(\operatorname{Tor}_2^R(k,k)\cong k\)。短正合列 \(0\to\mathfrak m\to R\to k\to0\) 遂给出 \(\operatorname{Tor}_1^R(\mathfrak m,k)\cong k\ne0\)。由 Tor 的平坦判据，\(\mathfrak m\) 不平坦；投射模必平坦，故它也不投射。

\ProseParagraphBreak
另一方面，在这条消解中，\(\Omega k=\operatorname{im}(R^2\to R)=\mathfrak m\)，而正合性给出
\[
\Omega^2k=\ker\bigl(R^2\xrightarrow{(x\;y)}R\bigr)
=\operatorname{im}\bigl(R\xrightarrow{\binom{-y}{x}}R^2\bigr)\cong R.
\]
最后的同构使用了最左边微分的单射性，所以 \(\Omega^2k\) 投射，而 \(\Omega k\) 不投射。这同时给出长度二的自由消解和阻止在第一步截短的一次 Tor。
\end{example}
